\chapter{DTI e DTA: Predizione delle interazioni Farmaco-Bersaglio}

\section{Il problema computazionale}

Una delle fasi più complesse della pipeline è la corretta predizione delle 
interazioni biochimiche tra il target (precedentemente identificato nella 
\textit{target prediction}) e i possibili farmaci.
Per introdurre modelli matematici e algoritmi di ML in questo step del 
processo è necessario ricondurre la ricerca a un problema computazionale 
ben definito, ovvero prevedere se, e con quale intensità, una determinata 
molecola interagirà con un target specifico.

Il nucleo di questo problema computazionale risiede nella capacità di 
modellare correttamente la \textbf{funzione di interazione} (o funzione 
di \textbf{scoring}).
\begin{equation}
    f(d,t) \xrightarrow{} y
\end{equation}
L'algoritmo ha il compito di filtrare lo spazio chimico come segue: 
data una proteina $t$ (\textit{target}), il modello esplora un dataset 
massivo di sequenze molecolari $\{d_1,d_2,...,d_n\}$ per identificare 
quelle che massimizzano la \textbf{funzione di scoring}.
Questi valori calcolati permettono di creare un \textit{ranking} dei 
potenziali composti candidati e delle interazioni più potenzialmente 
efficaci.
Questo processo è il passaggio cruciale della ricerca farmacologica 
moderna da problema sperimentale biochimico a problema di 
\textbf{ottimizzazione} e ricerca in uno \textbf{spazio latente} a 
dimensionalità alta.

La funzione può determinare, in base alla sua tipologia di output, due 
approcci principali:
\begin{itemize}
    \item \textbf{Classificazione Binaria} \{0,1\}, se il modello 
    fornisce un output booleano, che indica la assenza/presenza di 
    interazione con il target \cite{mondalAIdrivenDrugTarget2026a}.
    \item \textbf{Regressione}, se il modello ritorna un valore continuo 
    che quantifica la forza di legame, ovvero la 
    \textbf{Drug-Target Affinity (DTA)}.
\end{itemize}

\subsection{DTI: classificazione binaria}

\subsection{DTA: regressione del livello di affinità}

L'obiettivo principale della \textit{drug-target affinity} è la stima 
dell'energia libera di legame (ovvero l'affinità di legame) prevista tra 
farmaco e proteina-bersaglio \cite{mondalAIdrivenDrugTarget2026a}.
Questa stima include la previsione di parametri biologici come la 
\textbf{costante di dissociazione} ($K_d$), la \textbf{costante di 
inibizione} ($K_i$) o la \textbf{concentrazione inibente} ($IC_{50}$). 
Una volta calcolati, questi valori vengono convertiti in scala logaritmica 
per stabilizzare la convergenza durante la fase di \textit{training} dei 
modelli di Deep Learning \cite{kumarPredictingDrugtargetInteractions2026a}.

\subsection{Matrice di interazione e DTI Network}
Sia $D$ lo spazio chimico dei farmaci, e $T$ lo spazio biologico dei 
target, l'insieme delle relazioni tra i due può essere formalizzato 
attraverso una \textbf{matrice di interazione} farmaco-bersaglio 
(DTI matrix) \cite{mondalAIdrivenDrugTarget2026a}, in cui le righe 
rappresentano i ligandi, mentre le colonne rappresentano i bersagli 
biologici. 
Di conseguenza, ogni cella $M_{i,j}$ contiene il valore dell'interazione 
tra il farmaco $i$-esimo e il target $j$-esimo (in formato binario per la 
DTI e come valore continuo nella DTA).

Questo tipo di rappresentazione presenta tuttavia una problematica 
strutturale: le matrici di interazione sono infatti estremamente 
\textbf{sparse}, a causa del fatto che la maggior parte delle coppie 
$d$-$t$ non sono mai state testate sperimentalmente e non si hanno quindi 
informazioni a riguardo. 
Per superare tale limite, si ricorre a tecniche di 
\textit{Matrix Factorization} avanzate, come la \textit{Neighbourhood 
Regularized Logistic Matrix Factorization (NRLMF)}, che decompongono la 
matrice in componenti latenti per identificare pattern nascosti e 
prevedere interazioni sconosciute \cite{mondalAIdrivenDrugTarget2026a}.

Tuttavia, questo approccio tratta la matrice come un oggetto isolato e 
non sfrutta le relazioni strutturali tra i composti stessi, come la 
\textbf{similarità} chimica o funzionale tra molecole. 
Per integrare questo tipo di informazioni, le DTI matrix vengono 
reinterpretate come le matrici di adiacenza di una \textbf{DTI Network}.

Si tratta di un \textbf{modello di rete bipartita eterogenea}, in cui i 
nodi di una partizione (composti) sono collegati ai nodi dell'altra 
(target) tramite archi che rappresentano le interazioni note tra i due.
Questa riformulazione supera la staticità della matrice ``piatta'' e 
permette di modellare le relazioni intra-partizione, oltre a quelle 
farmaco-target, arricchendo ulteriormente il quadro informativo per la 
predizione \cite{mondalAIdrivenDrugTarget2026a}.

Modelli di ML supervisionato come il \textit{Bipartite Local Model} (BLM) 
sfruttano questa struttura per costruire dei classificatori \textit{locali} 
specifici per ogni nodo della rete.
A differenza di modelli globali che analizzano la rete nella sua interezza, 
il BLM addestra un classificatore separato per ogni farmaco e ogni target, 
rifacendosi alle interazioni nel \textit{neighbourhood} del nodo come 
segnale supervisionato. Questa strategia permette di catturare pattern 
specifici dei protagonisti dell'interazione, che i modelli globali 
generalizzerebbero eccessivamente.


\section{Metodologie di Predizione}
L'obiettivo è la costruzione di un modello che, in seguito ad un 
addestramento, sia in grado di fare generalizzazioni su coppie 
farmaco-target mai osservate prima, stimando correttamente la probabilità 
di interazione e la sua intensità.

Queste metodologie si differenziano in base al tipo di informazione su 
cui si basano:
\begin{itemize}
    \item I metodi \textbf{ligand-based}, tra cui rientrano i modelli 
    \textbf{QSAR (\textit{Quantitative Structure-Activity Relationship})}, 
    che associano le proprietà strutturali di un determinato composto con 
    la sua attività biologica.
    Storicamente implementati tramite modelli lineari, i modelli QSAR si 
    sono nel tempo evoluti verso modelli di ML e DL di complessità ed 
    efficienza maggiore, rimanendo comunque di natura \textit{ligand-based}.
    Anche la \textbf{modellazione farmacoforica} fa parte di questa 
    categoria, e serve ad identificare le caratteristiche strutturali 
    minime di un composto per ottenere l'interazione desiderata con il 
    bersaglio, stringendo il cerchio di possibili candidati (risulta 
    particolarmente utile nei casi in cui la struttura 3D della proteina 
    non è ricavabile \cite{mondalAIdrivenDrugTarget2026a}).
    \item I metodi \textbf{structure-based} che, al contrario dei 
    precedenti, richiedono la conoscenza della struttura 3D della proteina 
    target, al fine di simulare l'interazione chimica. Il rappresentante 
    principale di questa categoria è il \textbf{docking molecolare}, 
    che verrà trattato nella sezione successiva (\cref{subsec:docking}).
    \item I metodi \textbf{network-based}, comprendono la DTI network e 
    il modello BLM, già citati in precedenza. Questi rappresentano i 
    farmaci e i target attraverso i nodi di un grafo e sfruttano la 
    topologia della rete per analizzare nuove connessioni/interazioni.
    \item I metodi basati su \textbf{Machine Learning} e \textbf{Deep 
    Learning}, infine, rappresentano lo stato dell'arte.
    I modelli di ML classico apprendono pattern complessi a partire da 
    descrittori progettati manualmente, mentre le architetture di DL 
    estraggono automaticamente le informazioni facendo affidamento solo 
    al dataset di input, senza richiedere \textit{feature engineering}.
    Entrambi verranno trattati in dettaglio più avanti.
\end{itemize}

\subsection{Metodi ligand-based e structure-based}
\subsubsection{QSAR e modelli farmacoforici}

\subsubsection{Docking molecolare}
\label{subsec:docking}

\subsection{Machine Learning Classico}

Le metodologie computazionali per la predizione di interazioni 
farmaco-target hanno seguito un processo evolutivo motivato dalle 
limitazioni riportate dalle generazioni precedenti.
I primi approcci si basavano su descrittori delle proprietà molecolari 
dei composti progettati manualmente e basati unicamente sulle competenze 
e conoscenze scientifiche degli esperti di laboratorio.

\subsubsection{Modelli lineari}

I modelli di ML tradizionale si basano spesso sulla \textit{feature 
engineering} manuale, eseguita prima dell'addestramento del modello da 
ricercatori incaricati di selezionare le caratteristiche descrittive delle 
molecole, ovvero le loro proprietà fisico-chimiche.
Tra questi approcci si distinguono i \textbf{Modelli Lineari} 
\cite{odugbemiArtificialIntelligenceAntidiabetic2024}:
\begin{itemize}
    \item \textbf{MLR (Multiple Linear Regression)}
    \item La regressione dei \textbf{minimi quadrati parziali (PLS)}
    \item La \textbf{Regressione Lineare Bayesiana (RLB)}
\end{itemize}
Questi rappresentano tecniche statistiche in grado di generare una 
funzione lineare che descrive il legame tra i descrittori molecolari 
(le variabili $x$, dette ``indipendenti'') e l'esito della predizione, 
ovvero l'affinità di legame (variabili $y$, ``dipendenti'') 
\cite{odugbemiArtificialIntelligenceAntidiabetic2024}. 
In questo modo, il modello è capace di calcolare come i cambiamenti nelle 
proprietà dei composti influenzino l'attività biologica di essi. Diventa 
quindi possibile prevedere il comportamento di nuovi composti da analizzare, 
semplicemente ``consultando'' la linea di regressione e inserendo i nuovi 
parametri nella funzione.

Il limite di questi modelli, tuttavia, risiede proprio nella loro 
possibilità di catturare unicamente relazioni lineari, escludendo quindi 
pattern di interazione più complessi.

\subsubsection{Alberi decisionali e Random Forest}

Per superare questo ostacolo di complessità subentrano metodi più avanzati, 
come \textbf{Random Forest} e gli \textbf{alberi di decisione} (\textbf{DT}), 
che sono in grado di ricavare dai dati relazioni complesse, rappresentabili 
da funzioni non-lineari.

I \textbf{DT} sono una struttura dati ad albero in cui i dati vengono 
ricorsivamente partizionati in sottoinsiemi in base a caratteristiche 
specifiche come, in questo caso, particolari proprietà molecolari. 
In questo modo i rami della struttura rappresentano i possibili esiti del 
problema di classificazione/regressione (in base al tipo di task che viene 
affrontato, DTI o DTA).
Una volta ``riempito'', il DT può essere usato come modello predittivo, 
in grado di classificare nuovi composti in base al loro percorso guidato 
dalla radice fino alle foglie dell'albero 
\cite{odugbemiArtificialIntelligenceAntidiabetic2024}.
Queste tecniche sono efficaci e facili da interpretare, ma soffrono di 
\textbf{overfitting}: 
indeed, questi modelli tendono a ``memorizzare'' le caratteristiche dei 
dati del \textit{training set}, senza realmente apprendere dei pattern 
per generalizzare il comportamento dei vari input. In questo modo, le 
prestazioni risultano eccellenti durante la fase di training ma abbastanza 
scarse durante la fase di test o l'identificazione di nuovi dati.

La \textbf{Random Forest (RF)} è una tecnica di \textit{ensemble learning} 
progettata per ridurre il fenomeno di overfitting. È composta da più DT 
combinati e addestrati su sottoinsiemi differenti del dataset generale. 
La decisione finale avviene tramite voto di maggioranza, nel caso della 
DTI, o tramite media, nel caso della DTA.

\subsubsection{Support Vector Machines}

Oltre agli algoritmi \textit{tree-based}, sono utilizzate anche tecniche 
come \textbf{Support Vector Machines (SVM)}. Queste affrontano la 
classificazione DTI cercando, all'interno dello spazio dei descrittori, 
l'\textbf{iperpiano di separazione} con margine massimo tra coppie in 
interazione e non interagenti.
L'iperpiano ottimale è quello che massimizza la distanza tra i punti più 
vicini di ciascuna delle due classi, infatti più il margine è ampio e 
migliore sarà la capacità del modello di fare previsioni su coppie mai 
osservate. 
Questi punti ``di confine'' vengono chiamati \textbf{vettori di supporto} 
e garantiscono la stabilità del sistema 
\cite{montesinoslopezSupportVectorMachines2022}. 

Per gestire la non-linearità delle proprietà biologiche, SVM ricorre al 
\textbf{kernel trick}, che ``solleva'' i dati di una dimensione, 
garantendo sempre la separazione lineare. 
Il kernel non necessita delle coordinate dei punti nello spazio, ma 
ricava direttamente il loro \textbf{prodotto scalare} calcolandolo in uno 
spazio a dimensionalità superiore 
\cite{montesinoslopezSupportVectorMachines2022}. 

Il limite strutturale comune a tutti gli approcci di ML tradizionale è 
la dipendenza della predizione dalla qualità dei descrittori selezionati 
\cite{ahmadMachineLearningDrugtarget2026}.
L'inadeguatezza delle proprietà molecolari in esame può compromettere il 
lavoro del modello indipendentemente dalla sua potenza di calcolo o dalla 
sua efficienza.
Questo vincolo è proprio la motivazione del passaggio verso le architetture 
di \textbf{Deep Learning}, in grado di apprendere in modo autonomo le 
informazioni dai dataset.

\subsection{Deep Learning}
Il subentro del \textbf{Deep Learning} ha affrontato i problemi di DTI e 
DTA come un \textbf{task di apprendimento supervisionato}: il modello 
apprende la funzione
\begin{equation}
    f(d,t) \xrightarrow{} y
\end{equation}
in grado di generalizzare su coppie drug-target non osservate, 
riconoscendo pattern chimici e biologici di alta complessità.
Invece di ricevere parametri calcolati manualmente, questi modelli ricavano 
le caratteristiche più informative direttamente dai dati ``grezzi'', siano 
questi stringhe SMILES, sequenze amminoacidiche o grafi molecolari 
\cite{wangUnifiedSurveyDrugtarget2026a}.

\subsubsection{Modelli sequence-based (CNN, RNN/LSTM)}

I modelli \textit{sequence-based} affrontano il problema DTI trattando 
direttamente le stringhe SMILES o i fingerprint del farmaco e le sequenze 
delle proteine come dati testuali unidimensionali, senza richiedere input 
formalizzati ``su misura'' \cite{tangTCDTAPredictingDrugTarget2024a}.

\paragraph{CNN:}

Le \textbf{CNN} o \textbf{Reti Neurali Convoluzionali} rappresentano una 
delle prime architetture DL applicate con successo a questo tipo di input 
biologico.
Nascono principalmente per processare dati in formato a griglia come i 
pixel nelle immagini, ma possono essere adottate in ambito di predizione 
DTI e DTA per lavorare su task di QSAR, tramite \textbf{convoluzione 
unidimensionale} \cite{odugbemiArtificialIntelligenceAntidiabetic2024}.
Il processo si divide in tre fasi principali che seguono i diversi strati 
della rete:
\begin{itemize}
    \item \textbf{Layers Convoluzionali}, che scansionano le sequenze dei 
    composti e ricercano al loro interno dei pattern strutturali locali, 
    detti \textit{\textbf{motif}}.
    Il risultato del passaggio attraverso questi strati è la composizione 
    di una matrice \textbf{feature map}, che indica le feature estratte e 
    il grado di certezza nella loro rilevazione.
    \item \textbf{Layers di Pooling}, che riducono la dimensionalità delle 
    feature map ottenute, riassumendo solo le informazioni più importanti 
    e utili per la previsione di nuovi composti. Questo step, oltre a 
    rendere il modello più leggero, è fondamentale per la riduzione del 
    \textit{noise}, ovvero informazioni superflue che potrebbero 
    compromettere il risultato del procedimento con piccole variazioni di 
    misurazione e impurità nei dati.
    \item \textbf{Layers Fully Connected}, infine, sono strati densi che 
    combinano tra loro le caratteristiche apprese dalle mappature nei 
    livelli precedenti, per generare una previsione finale.
\end{itemize}

Questo tipo di reti neurali sono facilmente scalabili e molto efficienti 
ma, facendo uso di informazioni 1D, fanno anch'esse fatica ad incorporare 
strutture tridimensionali più complesse e a catturare le dipendenze 
biochimiche tra elementi \cite{houDGCADTADeepGraph2025}.

\paragraph{RNN:}

Le \textbf{RNN (Recurrent Neural Networks)} sono un'altra architettura 
di Deep Learning progettata per l'elaborazione di dati sequenziali. 
A differenza delle CNN, però, queste reti non sono in grado di modellare 
in modo corretto e affidabile le dipendenze tra atomi tra loro distanti 
\cite{wangUnifiedSurveyDrugtarget2026a}.

Operano mantenendo un \textbf{layer nascosto}, che si evolve con 
l'elaborazione della sequenza di input. 
Nel contesto della predizione di interazioni con il target farmacologico, 
le RNN, come le CNN, analizzano efficacemente le stringhe SMILES dei 
composti per apprendere i \textbf{motifs} strutturali e i pattern nelle 
relazioni che determinano l'affinità di legame.

\subsubsection{Modelli graph-based (GNN)}

I modelli \textit{sequence-based}, anche se risultano più semplici nella 
formattazione degli input e più facilmente scalabili, faticano a prendere 
conoscenza del contesto tridimensionale del sito di legame 
\cite{wangUnifiedSurveyDrugtarget2026a}. Questa perdita di informazioni 
può essere superata grazie alla rappresentazione a grafo tipica dei modelli 
trattati in questa sezione.

Come accennato brevemente nel capitolo precedente, questo tipo di rappresentazione 
permette l'elaborazione dei contesti da parte delle \textbf{GNN (Graph 
Neural Networks)}, che operano simulando la topologia e le relazioni 
spaziali reali dei composti chimici (legami tra atomi \textit{intra} e 
\textit{inter} composto) propagando informazioni tra i nodi tramite il 
meccanismo di \textit{message passing}. Questo processo avviene 
in due passaggi iterativi \cite{wangGraphNeuralNetworks2023}:
\begin{itemize}
    \item \textbf{Computazione e Aggregazione:} Il modello calcola un 
    messaggio elaborando le caratteristiche del nodo centrale, con quelle 
    dei nodi e degli archi ad esso adiacenti.
    \item \textbf{Aggiornamento:} Lo stato del nodo viene aggiornato 
    integrando il nuovo messaggio con la sua informazione precedente.
\end{itemize}

Queste informazioni elaborate vengono poi lette tramite l'operazione di 
\textbf{Readout}. Attraverso un meccanismo simile al pooling, le 
caratteristiche di tutti i nodi vengono aggregate e aggiornate in un 
unico vettore di caratteristiche globali, in grado di rappresentare lo 
stato dell'intera molecola \cite{wangGraphNeuralNetworks2023}.

In generale, nel contesto di drug discovery le GNN stanno rivoluzionando 
il concetto di tradizionale \textbf{Computer-Aided Drug Design} (CADD), 
trasformando il processo in una ricerca automatizzata.
Queste architetture consentono di esplorare lo spazio chimico per generare 
e ottimizzare (come nella \textit{Lead Optimization}) nuove entità 
molecolari, apportando migliorie nelle proprietà desiderate, come 
bioattività, stabilità e sintetizzabilità. 
Queste permettono anche uno studio precoce di proprietà ADMET e della tossicità, 
portando a una riduzione drastica di fallimenti nei test clinici 
\cite{wangGraphNeuralNetworks2025}.

Nel contesto del problema computazionale DTI/DTA, le GNN risultano fondamentali nello screening virtuale per prevedere con efficienza le affinità di legame tra farmaci e target attraverso l'implementazione di architetture a doppia branca (\textit{dual-branch}). In questo paradigma, il problema predittivo viene convertito in un compito di fusione di informazioni estratte in parallelo da due grafi distinti:

\begin{itemize}
    \item \textbf{Branca del Farmaco ($d$):} La piccola molecola viene convertita in un grafo molecolare in cui gli atomi rappresentano i nodi (caratterizzati da proprietà chimiche come numero atomico, valenza, ecc...) e i legami chimici costituiscono gli archi. Algoritmi come le \textit{Graph Convolutional Networks} (GCN) scansionano questa struttura per estrarre pattern locali e stereochimici stabili.
    \item \textbf{Branca del Target ($t$):} La proteina viene modellata come un grafo di contatto dei residui. In questa struttura, i nodi rappresentano i singoli amminoacidi, mentre gli archi vengono tracciati tra nodi la cui distanza euclidea nello spazio 3D è inferiore a una determinata soglia, preservando la topologia del sito di legame \cite{ahmadMachineLearningDrugtarget2026}.
\end{itemize}

Una volta che i messaggi sono stati propagati all'interno delle rispettive strutture, le operazioni di \textit{readout} generano \textbf{due vettori di feature ad alta dimensionalità}. Questi vettori vengono successivamente combinati all'interno di un modulo di interazione, spesso arricchito da meccanismi di attenzione (\textit{Graph Attention Networks} - GAT) per pesare l'importanza relativa dei vari atomi e residui nel complesso farmaco-bersaglio \cite{ozsariDTAGNNToolkitConstructing2026}. Il vettore risultante viene infine processato da strati \textit{Fully Connected} per mappare l'interazione nell'output finale $y$, traducendosi in una classificazione booleana per i task di DTI o in un valore continuo di affinità ($K_d$, $K_i$, $IC_{50}$) per i task di DTA.

\subsubsection{Modelli attention-based e Transformer}

Le architetture basate sul meccanismo di \textbf{attenzione} e sui 
\textbf{Transformer} rappresentano l'attuale stato dell'arte nella 
predizione di DTI e DTA. Il loro punto di forza principale risiede nella 
capacità di catturare \textbf{dipendenze a lungo raggio} all'interno delle 
sequenze molecolari e proteiche \cite{ahmadMachineLearningDrugtarget2026} \cite{mondalAIdrivenDrugTarget2026a}, superando il limite strutturale delle CNN 
e delle RNN, che tendono a perdere il contesto globale quando le sequenze 
diventano molto lunghe.

\paragraph{Il meccanismo di attenzione}

L'\textit{attention} si basa sulla premessa che non tutte le parti di 
una sequenza sono ugualmente rilevanti per una previsione. 
Si rende quindi necessario \textbf{pesare selettivamente} le sottosequenze più 
informative, ad esempio, applicato alle interazioni farmaco-bersaglio, i residui amminoacidici del sito di legame, ignorando le regioni che non contribuiscono all'interazione.
% figura? heatmap


\paragraph{Architetture Transformer per DTI e DTA}

I \textbf{Transformer} estendono il meccanismo di attention trattando le stringhe SMILES e le sequenze 
amminoacidiche (FASTA) come se fossero un linguaggio naturale da 
interpretare. 
A differenza delle RNN, che elaborano i token uno alla 
volta in modo sequenziale, i Transformer processano l'intera sequenza 
simultaneamente, rendendo il training più efficiente e la cattura del 
contesto globale più affidabile \cite{wangUnifiedSurveyDrugtarget2026a}.
% figura di architettura generale di un Transformer encoder

Tra i modelli più rappresentativi si distinguono:
\begin{itemize}
    \item \textbf{MolTrans} e \textbf{MT-DTI}, che applicano encoder 
    Transformer sia al farmaco che alla proteina per apprendere 
    rappresentazioni consapevoli del contesto. Questi modelli mostrano 
    ottime prestazioni anche in scenari \textit{cold-start}, ovvero in 
    presenza di coppie farmaco-target nuove mai osservate durante il training \cite{mondalAIdrivenDrugTarget2026a}.
    \item \textbf{TC-DTA}, un modello ibrido che adotta un encoder 
    Transformer per la sequenza proteica e una CNN per la rappresentazione 
    del farmaco, ottimizzando l'estrazione di informazioni semantiche 
    dagli amminoacidi pur mantenendo l'efficienza computazionale della rete
    convoluzione per il lato molecolare \cite{tangTCDTAPredictingDrugTarget2024a}.
    \item \textbf{FusionDTI}, che integra token strutturali proteici e 
    stringhe SELFIES tramite Transformer, con l'obiettivo di ottenere 
    predizioni a grana fine sul sito di legame specifico.
\end{itemize}

\paragraph{Graph Transformer}
Un'ulteriore evoluzione consiste nell'applicare la logica dei Transformer 
direttamente alle strutture a grafo delle molecole, dando origine ai 
cosiddetti \textbf{Graph Transformer}. Questi modelli incorporano nei 
meccanismi di \textit{attention} informazioni strutturali specifiche del grafo, 
come la \textbf{centralità dei nodi} e le \textbf{distanze spaziali} tra 
atomi, che le GNN standard non sfruttano in modo esplicito.
Il modello \textbf{GraphormerDTI} rappresenta un esempio di questa 
categoria applicata al problema DTI, dimostrando miglioramenti rispetto 
alle GNN classiche grazie all'integrazione di queste feature informative strutturali 
nell'attenzione \cite{kumarPredictingDrugtargetInteractions2026a}.

\paragraph{Modelli linguistici pre-addestrati (PLM)}
I Transformer costituiscono la spina dorsale dei cosiddetti 
\textbf{Foundation Models}, ovvero modelli pre-addestrati su database 
massivi in modalità auto-supervisionata e successivamente adattati a 
compiti specifici tramite \textit{fine-tuning}. Questo paradigma 
\textit{pre-train / fine-tune} permette di sfruttare la conoscenza 
biochimica e chimica generale acquisita su larga scala verso i task di DTI e DTA 
con dataset di dimensioni più ridotte \cite{wangUnifiedSurveyDrugtarget2026a}.

\paragraph{Vantaggi e limiti}
Uno dei punti di forza più rilevanti di questi modelli in ambito 
farmacologico e/o medico è la loro \textbf{interpretabilità}: i pesi di attenzione 
possono essere visualizzati attraverso delle \textbf{mappe di attenzione}, che indicano 
quali atomi o residui hanno guidato maggiormente la predizione \cite{kumarPredictingDrugtargetInteractions2026a}. Questo 
strumento offre ai ricercatori la possibilità di \textbf{validare 
biologicamente} il risultato computazionale, verificando se le regioni 
evidenziate corrispondono a siti di legame noti.
% figura saliency map ?

Il principale limite rimane la \textbf{complessità computazionale}: 
l'\textit{attention} ha complessità quadratica rispetto alla lunghezza della 
sequenza, il che rende il training particolarmente oneroso in termini di 
tempo e memoria quando si lavora con sequenze proteiche molto lunghe \cite{mondalAIdrivenDrugTarget2026a}

%\subsubsection{Modelli ibridi}